%!TEX program = uplatex
\RequirePackage{plautopatch}
\documentclass[uplatex,dvipdfmx]{jlreq}
\usepackage{amsmath,amssymb}

\begin{document}

\noindent
{\large\bfseries モーデル・ヴェイユ格子　入門}

\medskip
\noindent
{\large\bfseries \hfill 塩田 徹治（立教大学・理）}

\bigskip\bigskip

１９８９年春にはじまるモーデル・ヴェイユ格子（と名付けたもの）との出会いは、
私にとっても文字通り新鮮な「数学との遭遇」であった。
旧知のひとに出会い、昔よりはるかにそのひとのことをよく理解する機会にめぐりあったようなものだ。
当時、一つの楕円曲線の有理点を具体的に求めることをコンピュータで実験してみた所、
特別な形の２４０個の有理点が得られた。
この結果が、ランク８のルート格子 $E_8$ の２４０個の最短ベクトルの決定に
相当することに気付いた時が、物事の転回点となった。

\bigskip

モーデル・ヴェイユの定理によれば、有理数などを係数とする楕円曲線
（やアーベル多様体）の有理点の全体は、
有限生成なアーベル群となる；それをモーデル・ヴェイユ群とよぶ。
モーデル・ヴェイユ格子（MWL）の理論の基本的アイデアは、
「モーデル・ヴェイユ群に、適当な内積をいれて、格子としてみる」ということである。
この立場からみると、予想外に豊富な事柄が自然なつながりをもって見えてくるのである。
ある人は私のMWLの話を評して、箱庭のようだ、といったことがある。
代数、幾何、数論、代数幾何、トポロジー、など色々なものを、
木々や草花、池や水の流れ、丘や築山などにたとえたのだろうか。

\bigskip

モーデル・ヴェイユ格子に関して、できるだけ分かりやすい紹介を、とのご要望をうけて、
３回の講演の予定を、以下のようにたててみた。

\bigskip\bigskip

\noindent\textbf{I\@.　モーデル・ヴェイユ格子の背景から定義まで}
\begin{enumerate}
  \setlength{\itemsep}{0pt}
  \item 格子と Sphere Packing（球のつめこみ）
  \item 楕円関数、楕円曲線、モーデル・ヴェイユの定理
  \item 楕円曲面上の交点理論
  \item モーデル・ヴェイユ格子（MWL）
\end{enumerate}

\bigskip

\noindent\textbf{II\@.　$E_6, E_7, E_8$ の背後にひそむもの}
\begin{enumerate}
  \setlength{\itemsep}{0pt}
  \item ルート格子　$E_6, E_7, E_8$　とその双対格子
  \item 有理楕円曲面のMWL
  \item MWLから生ずる代数方程式とガロア表現
  \item 特異点の変形とMWL
  \item 古典再訪：３次曲面上の２７本の直線、４次曲線の２８本の双接線
\end{enumerate}

\bigskip

\noindent\textbf{III\@.　応用と展望}
\begin{enumerate}
  \setlength{\itemsep}{0pt}
  \item Sphere Packing への応用（標数 $p$ の超特異現象の意外な効用）
  \item ランクの高い楕円曲線の構成
  \item より一般の枠組みでのMWL展望
\end{enumerate}

\bigskip\bigskip

講演を聞かれた方に、なるほど面白いと感じて頂けるならば幸いである。

\bigskip

\noindent
（お手近な参考文献：雑誌「数学」43巻（1991）所収の論説　「MWLの理論と応用」）

\end{document}